Mixed human-AI work requires more than generation and tool calling. It requires a durable thread that can survive ambiguity at intake, preserve commercial boundaries, route to the right form of supply, and demand proof appropriate to the execution mode. This need becomes most obvious when the work includes embodied steps. Physical presence, local access, inspection, and witnessed completion cannot be treated as optional just because they are harder for digital planners to execute directly.

We have argued that a request-rooted orchestration model offers one way to address this gap. By keeping \texttt{Request} as the durable root, separating \texttt{Commitment} from \texttt{Fulfillment}, preserving proof through \texttt{Artifact}, and recording business-meaningful history through \texttt{RequestEvent}, the system can keep the work thread intact across both digital and embodied execution. The lead-first orchestration rule further helps resist task explosion and tool-shaped underplanning.

The broader point is not that every request should be solved by AI, nor that every planner failure can be fixed by better prompting. The point is that the surrounding work model determines what kinds of failure become visible or invisible. If the system has no durable place for embodied constraints and proof obligations, it will tend to produce false closure. The deterministic benchmark in this artifact makes that failure observable: request-rooted planning preserves full embodied recall and zero false completion on the current four-scenario pack, while task-first and direct-tool baselines fail for different reasons, namely over-decomposition and digital-only collapse. The live benchmark then shows the complementary real-model picture: current models often recover the correct lead and remain action-safe, but they still struggle to emit the full canonical proof-bearing structure that Boreal requires.

A request-rooted model therefore offers a stronger basis for trustworthy mixed human-AI fulfillment. The next step is not merely larger language models. It is broader evaluation on live traces, more diverse embodied scenarios, stronger linkage between planning outputs, delivered artifacts, and settlement-safe business truth, and narrower structured interfaces that reduce the gap between semantic planning competence and exact durable contract emission.
